% =========================================================
% Proof Stub: Dirichlet Test
% Source: volume-iii/analysis/sequences/notes/series/notes-series-tests.tex
% =========================================================

% *** HARDER PROOF — attempt after foundational results settled ***

\clearpage
\phantomsection
\hypertarget{proof-RA-ABB-T-dirichlet-test}{}

\subsubsection[Dirichlet Test]{Proof --- Dirichlet Test}

\bigskip

\noindent
\textbf{Source.} Abbott, \emph{Understanding Analysis}; see \texttt{notes-series-tests.tex}.

\vspace{0.75em}

\noindent
\textbf{Goal.} If $(a_n)$ is decreasing to 0 and the partial sums of $(b_n)$ are bounded, then $\sum a_n b_n$ converges.

\vspace{0.75em}

\noindent
\textbf{Logical form.}
\[
  a_n \searrow 0,\; |B_n| \le M \Rightarrow \sum a_n b_n \text{ converges}.
\]

\vspace{0.75em}

\noindent
\textbf{Proof strategy.} Use Abel's summation by parts: $\sum_{k=m}^n a_k b_k = a_n B_n - a_m B_{m-1} - \sum_{k=m}^{n-1}(a_{k+1}-a_k)B_k$. Bound using $|B_k| \le M$ and $a_n \to 0$.

\vspace{0.75em}

\noindent
\textbf{Proof.}
\begin{proof}
\Claim If $(a_n)$ is decreasing to 0 and the partial sums of $(b_n)$ are bounded, then $\sum a_n b_n$ converges.

\Given 

\Goal 

\Strategy Use Abel's summation by parts: $\sum_{k=m}^n a_k b_k = a_n B_n - a_m B_{m-1} - \sum_{k=m}^{n-1}(a_{k+1}-a_k)B_k$. Bound using $|B_k| \le M$ and $a_n \to 0$.

\medskip

% --- summation by parts formula ---
% --- bound each term using |B_k| ≤ M ---
% --- a_n → 0 ⟹ a_n B_n → 0 ---
% --- ∑(a_k - a_{k+1}) telescopes to a_m; bounded ---

\end{proof}

\begin{remark*}[Proof shape]
Summation by parts (Abel summation) converts the problem into one involving the bounded partial sums of $(b_n)$ and the telescoping of $(a_n)$.
\end{remark*}

\begin{remark*}[Dependencies]
\begin{itemize}
  \item Abel summation by parts formula.
  \item $(a_n)$ decreasing to 0 $\Rightarrow$ total variation $\sum |a_{k+1} - a_k| = a_1 - \lim a_n = a_1$.
  \item Partial sums of $(b_n)$ bounded by $M$.
\end{itemize}
\end{remark*}
